\section{已废止的紧凑与 Nielsen 路线}\label{app:retired-routes}

本附录记录在第~\ref{sec:johnson-presentation} 节采用 Johnson 替换之前曾检验并否决的路线。
此处材料均不用于正文论证。

符号 AR 公式与标准穿孔圆盘 braid 本身不能确定约化 handlebody 边界中嵌入的 $44$-passage 领圈。
显式 19 步 Nielsen 代表约化 $m_2$ 长度为 $309$、42 个 detector 通道，而紧凑代表长度为 $311$、44 个通道；因此它不是公开领圈词的来源。

对 $x^{-1}$ 的内共轭恢复 44 个带基点 passage 并通过 $F_3$ 的自由基检验，但不改变无基点循环脊类；它不是嵌入领圈见证。
对非内 IA 修正的有界搜索找到改变 detector 的 44 通道候选，但没有满足互补 handlebody 检验的符号置换 dual-meridian 扩张。

紧凑直脊词在 $F_3$ 上单射但不满射，故它们不能是保持 handlebody 的 $\psi_A\in\operatorname{Aut}(F_3)$ 的像。
正文中的 Johnson $\alpha_{ij}$ 分解规避了此阻碍。

在词层面，309 与 311 字母代表收集到同一正规形 $y^{40}z^{269}$，但所需的 framing 运输依赖于 $\operatorname{lk}(m_2,r_{yz})$，而词本身不能确定它。
为 Johnson CS 柄图构造的 railroad 约化 PD（引理~\ref{lem:P0d-link}）给出 $\operatorname{lk}(m_2,r_{yz})=0$。
